\scanpage{112}
\uline{Lemme 1} \inserted{(cas $\unsure{c|S}$ fini $\neq1,2$)} Soient $s,t\in\widetilde E_a$, alors $\exists\,s_1\in\widetilde E_a$, et $u,v\in\widetilde E_b$ tels que $u$ soit lié à $s$ et $s_1$, $v$ lié à $s_1$ et $t$.

% slanted
\margin{C'est \uline{faux} si $c=1,2$}

Soit $\widetilde E_{b}(s)=\{u\in\widetilde E_b\mid u\text{ lié à }s\}$, de $=$ $\widetilde E_b(t)$, ils sont de cardinal $c$ (des sections de $\widetilde E_b$ sur $E_b$) Si $\widetilde E_b(s)\cap\widetilde E_b(t)\neq\emptyset$, il suffit de prendre $u=v$ \deleted{\unsure{dans|élt}} dans cette intersection et $\bar s_1=s$. Sinon, comme $\operatorname{Card}\widetilde E_b=2c$, on a

\[ \widetilde E_b\deleted{\unsure{\varnothing}}=\underset{\substack{\shortparallel\\ B'}}{\widetilde E_b(s)}\amalg\underset{\substack{\shortparallel\\ B''}}{\widetilde E_b(t)} \]

(\unsure{ou mm|ou non|ou même|comme même} $\widetilde E_b(t)=\widetilde E_b(s')$ et sans doute cela impliquera $t=s'$, mais pour l'instant on ne peut le prouver). On cherche \deleted{\unsure{donc}} $\bar s_1$ tel que $\widetilde E_b(\bar s_1)\cap\widetilde E_b(s)\neq\emptyset$, $E_b(\bar s_1)\cap\widetilde E_b(t)\neq\emptyset$.

Si un tel $\bar s$ n'existait pas, on aurait $\forall\,\unsure{r|\bar r}\in\widetilde E_a,\ \widetilde E_b(\unsure{r|\bar r})$ est contenu dans \deleted{$\widetilde E_b(s)$} $B'$ ou $B''$ et pour des raisons de cardinal (comme il est lui $=$ de cardinal $c$) on aurait

\[ \forall r\in A=\widetilde E_a,\qquad \text{on a}\quad \widetilde E_b(r)=B'\ \text{ou}\ \widetilde E_b(r)=B'' \]

Soit

\[
\begin{aligned}
A'&=\{r\in A\mid\widetilde E_b(r)=B'\},& A''&=\{r\in A\mid\widetilde E_b(r)=B''\}\\
&&&=\complement_AA'
\end{aligned}
\]

donc \deleted{$A'=\sigma_{\underline t_0}A'',\ A''=\sigma_b A'$, et \unsure{pour}} $\forall$ couple $r,r'$ de pts de $A$ liés, un et un seul de $r,r'$ est dans $A'$, l'autre dans $A''$ -- donc $A',A''$ sont des sections complémentaires de $\widetilde E_a$ sur $E_a$.

Soit \deleted{$r_1,r_2\in\widetilde E_a$, \unsure{avec}} \deleted{\unsure{et}} $r_1,r_2\in A'$, \inserted{$r_1\neq r_2$} \deleted{$r_2\in A''$} \deleted{$r_2\neq r'_1$} \inserted{$r_1\in A'\ r'_2\in A''$} \deleted{de sorte que $\unsure{\nexists}$} et $r'_2=(r_2)'$ ; \deleted{il n'$\unsure{\exists}$} alors $\nexists\,u\in\widetilde{\unsure{E|B}}_b$ lié à $r_1,r'_2$ à la fois, donc l'ens des pts de $\widetilde E_c$ liés à $r_1$ et $r'_2$ \unsure{à|est} la fois est

% archive sheet number: 111
